\section{The research arc}
\label{sec:arc}

This project has run for six versions. The arc matters to the reader of this paper for one reason:
the v0.7 evaluation of \S\ref{sec:results} is only interpretable against what came before it, and
what came before it includes two occasions on which this project's own published claims turned out
to be artefacts of its measurement rather than properties of its policies. A reader who wants to
know why the apparatus of \S\ref{sec:methodology} is as heavy as it is will find the answer here.

\subsection{v0.2 to v0.4: building the object}
\label{sec:arc-early}

The first versions established the domain and the inference. A count-conditioned, empirically tuned
linear policy in a TypeScript simulator became a rules-exact C++ engine with an exact reference
posterior, a fast deployed approximation, a twenty-feature ask score and a value-based declaration
rule. The substantive contribution of that period is the object described in
\S\ref{sec:engine-belief}: an exactly computable observer-conditioned deal posterior in a game where
nobody had computed one, cheap enough to run in all six seats simultaneously.

That period also set the pattern this paper spends a great deal of effort undoing. Its evaluation
was a few thousand games per cell against a handful of opponents, its ablations were unpaired, and
its conclusions were drawn from point estimates. \S\ref{sec:corrections-corpus} lists what that cost.

\subsection{v0.5: a failure mode found, removed, and honestly reported as worth nothing}
\label{sec:arc-vfive}

v0.5 diagnosed a two-question deadlock in mirror play --- two policies alternately asking each other
for cards neither would surrender, ending in misdeclaration --- and removed it. It then reported, in
terms, that removing it was worth almost nothing in win rate.

That report is the honest ancestor of this one. It is also the point at which the project's central
difficulty became visible: a release that removes a real defect and gains no measurable strength is
either at the ceiling of its policy class or is measuring badly, and those two possibilities call
for opposite responses.

\subsection{v0.6: the measurement crisis}
\label{sec:arc-vsix}

The v0.6 cycle set out to find several points of accessible headroom that the accumulated design
register described. It found none of it, and in the process discovered that most of the apparatus
was broken. Four results from that cycle are load-bearing for this paper.

\subsubsection{Exchangeable ties are irreducible, and the tie-break is a coordination decision}
\label{sec:arc-vsix-ties}

At \vsixTieShareFive\% of \emph{contested} ask decisions two or more candidates score
\emph{bit-for-bit identically}, and the winner is whichever the enumerator emitted first. \vsixTieSameSuitFive\% of those ties are
two cards of one half-suit at one target, and the \emph{exact} posterior --- not merely the approximation the policy
deploys --- \emph{separates} them in \vsixTieExactSepPctFive\% of cases, i.e. assigns them identical
probability in every one. Every rule that ranks them by that probability therefore
realises the same hit rate, and so does enumeration order.

The largest item on the inherited register was not a defect. It was a property of the information
set.

What is \emph{not} irreducible is the dependence on enumeration order itself. An unstable sort order
deciding half of all contested decisions is a correctness problem rather than a strength problem,
and \S\ref{sec:domain-team}'s common-knowledge property says how to fix it: replace the order with a
hash of the public event stream, which every teammate can compute and which no seat's private hand
influences. That switch is one of the five keys of the configuration this paper evaluates
(\S\ref{sec:configuration}), and \S\ref{sec:rejected-tie} reports what happened when a later cycle
tried to do better than randomising inside the tie group.

\subsubsection{The exact posterior is a worse predictor than the approximation}
\label{sec:arc-vsix-belief}

Scored as predictors on identical states --- mean predictive log loss over unresolved cards, and the
realised hit rate of each posterior's own argmax ask --- the ranking is:

\begin{center}\small
\begin{tabular}{lrr}
\toprule
Posterior & Log loss (nats) & Argmax hit rate \\
\midrule
Exact, uniform deal prior & \vsixBeliefExactNLL{} & \vsixBeliefExactHit\% \\
Rules posterior, no policy prior & \vsixBeliefNoPriorNLL{} & \vsixBeliefNoPriorHit\% \\
\textbf{Deployed approximation} & \textbf{\vsixBeliefShippedNLL{}} & \textbf{\vsixBeliefShippedHit\%} \\
\bottomrule
\end{tabular}
\end{center}

The exact object is the worst row. The policy prior is the entire difference: exactness here is
bought by discarding behaviour, because the block construction of \S\ref{sec:engine-belief} has no
parameter through which behaviour could enter. Exactness and behaviour-awareness are not available
in the same object in this game.

This closed an experiment three studies had listed as outstanding, and it established the standing
rule this project has applied ever since --- ``exact Bayesian inference'' is an assumption to test.
\S\ref{sec:rejected-alloc} is the second time that rule paid, on a different channel.

\subsubsection{The optimiser's curse in determinized search}
\label{sec:arc-vsix-search}

v0.6 built a determinized search that samples full deals from the public posterior and plays each
forward with all six players reconstructed at their own information sets --- not clairvoyant, which
\S\ref{sec:domain-not} explains would be degenerate.

Chosen by an unguarded argmax over sampled means, that search is \textbf{\vsixCurseGap{}~points
worse} than the blueprint it searches from, measured against a control that forces the same code to
return the blueprint's own pick at identical budget, seed and sample size. The reason is the
statistic, not the search: one determinization's return is the final half-suit differential, whose
spread is an order of magnitude larger than genuine differences between candidates, so the argmax
selects on residual noise and does so more confidently the more candidates it is offered.

A paired lower-confidence-bound deviation rule recovers all of it. Guarded, the search is worth \vsixSearchPooled\% against \vfive{} over \vsixSearchPooledN{} games at
\vsixSearchPooledBanks{} banks. The transferable result is the guard rather than the search, and the
guarded search --- shipped switched off in v0.6 --- is one of the five keys of the configuration
this paper evaluates.

\subsubsection{The evaluation could not resolve a point, and the optimiser could not move a policy}
\label{sec:arc-vsix-apparatus}

Two apparatus failures, both of which this paper's methodology exists to prevent.

\textbf{The evaluation.} Pooling every same-pair cell in the v0.6 battery gave \vsixHeadPooled\%
with a naive $z$ of \vsixHeadZ{} --- a null, and an earlier draft of that report published it as
one. Re-run at high power across \vsixBigFiveBanks{} disjoint banks, the same comparison is
\vsixBigFive\% over \vsixBigFiveN{} games $[\vsixBigFiveLo, \vsixBigFiveHi]$, above parity on
\emph{all \vsixBigFiveAbove{}} banks, worst bank \vsixBigFiveMin\%. The two cells that had gone the
other way returned above parity at the larger sample. The effect was always there; the instrument
could not see it.

Five mechanisms in that cycle cleared a 95\% paired interval at the per-cell budget this project had
used since v0.3 and returned \emph{exactly zero} at three times it. That is the fact that makes
every sample size in \S\ref{sec:methodology-power} non-negotiable.

\textbf{The optimiser.} The previous fitting harness could not move a policy: one step size for
thirty-four coordinates spanning two and a half orders of magnitude, an unpaired objective, and an
objective documented as a worst case that was in fact a weighted mean. Its forty-generation refit
was indistinguishable from sampling noise, so the vector it ``found'' was the previous version's.
The repaired optimiser then found a policy that trades nothing head to head for a large replicated
gain against the deception archetype the project owner brought from live play
(\vsixPoolWithholderGain{}~points).

\subsection{What v0.6 handed forward}
\label{sec:arc-handoff}

A policy about nine tenths of a point stronger than its predecessor, a design register most of whose
entries had just been closed as worth nothing, and a diagnosis: \emph{at the budget this project had
always used, this policy class looks flat.} v0.6 stated plainly that it did not know whether the
class \emph{is} flat or whether the instrument was still too coarse.

v0.7 is the cycle built to answer that, and it was organised so the answer could not be argued into
being positive. Five phases ran under separate briefs, each starting from a cleared context:

\begin{enumerate}
\item \textbf{Rebuild the instrument} and characterise what it can and cannot see
      (\S\ref{sec:instrument}).
\item \textbf{Generate adversaries in the open} and find what actually beats the v0.6 frontier ---
      not the deployed policy, the frontier (\S\ref{sec:instrument-adversaries}).
\item \textbf{Build candidate architectures} against the register's ranked hypotheses
      (\S\ref{sec:rejected}).
\item \textbf{Bake off and freeze} exactly one configuration (\S\ref{sec:configuration}).
\item \textbf{Evaluate it} on sealed material under a protocol committed in advance
      (\S\ref{sec:protocol}, \S\ref{sec:results}).
\end{enumerate}

The rest of this paper is that programme and its result.
